\ea\label{ex:text6-1}
O koverau na pura nana rasuu.\\
\gll o	koverau	na	pura	nana	rasuu\\
     \textsc{art}3.\textsc{sg}	koverau.bamboo	\textsc{tm}2	grow	3\textsc{sg}.\textsc{ipfv}	jungle\\
\glt `The koverau bamboo grows in the jungle.'
\z

\ea\label{ex:text6-2}
O koverau saka tau pura vatotobin tamuana bata haana.\\
\gll o	koverau	saka	tau	pura	va-totobin	tamuana	bata	haa=na\\
     \textsc{art}3.\textsc{sg}	koverau.bamboo	\textsc{neg}1	\textsc{tm}5	grow	\textsc{advr}-straight	always	\textsc{cont}	\textsc{neg}2=3\textsc{sg}.\textsc{ipfv}\\
\glt `The koverau bamboo might not always grow straight.'
\z

\ea\label{ex:text6-3}
Tau hoa vanaavihu batana bara tau obete gunaha suku batana bona maa hum gunaha.\\
\gll tau	hoa	vanaavihu	bata=na	bara	tau	obete	gunaha	suku	bata=na	bona	maa	hum	gunaha\\
     \textsc{tm}5	arise	opposite.direction	\textsc{cont}=3\textsc{sg}.\textsc{ipfv}	and2	\textsc{tm}5	lie	go.down	follow	\textsc{cont}=3\textsc{sg}.\textsc{ipfv}	\textsc{obj}.\textsc{art}2.\textsc{sg}	\textsc{plm}	place	go.down\\
\glt `(It) might grow in the opposite direction and lie on the ground following downwards the slopes of a mountain.'
\z

\ea\label{ex:text6-4}
Ei be pura tea maa vaeereere.\\
\gll ei	be	pura	te=a	maa	vaereere\\
     \textsc{dem}2	\textsc{compl}2	grow	\textsc{prep}=\textsc{art}2.\textsc{sg}	\textsc{plm}	edge.of.mountain\\
\glt `This (means) that (it) grows on the edges of mountains,'
\z

\ea\label{ex:text6-5}
Pura vageegeno batari amaa meha.\\
\gll pura	vageegeno	bata=ri	a=maa	meha\\
     grow	crookedly	\textsc{cont}=3\textsc{pl}.\textsc{ipfv}	\textsc{art}2.\textsc{sg}=\textsc{plm}	other2.\textsc{sg}\\
\glt `Some grow crookedly.'
\z

\ea\label{ex:text6-6}
O koverau na vaarutaa ki nana bono kururu.\\
\gll o	koverau	na	vaa-rutaa	ki	nana	bono	kururu\\
     \textsc{art}3.\textsc{sg}	koverau.bamboo	\textsc{tm}2	\textsc{caus}-small	\textsc{dat}	3\textsc{sg}.\textsc{ipfv}	\textsc{obj}.\textsc{art}3.\textsc{sg}	kururu\\
\glt `The koverau bamboo is smaller than the kururu bamboo.'
\z

\ea\label{ex:text6-7}
O naono o mataa, na paku raara bono tahitahi.\\
\gll o	naono	o	mataa	na	paku	ra=ara	bono	tahi{$\sim$}tahi\\
     \textsc{art}3.\textsc{sg}	plant	\textsc{art}3.\textsc{sg}	good	\textsc{tm}2	make	1\textsc{pl}.\textsc{in}.\textsc{ipfv}=1\textsc{pl}.\textsc{in}.\textsc{pron}	\textsc{obj}.\textsc{art}3.\textsc{sg}	\textsc{red}{$\sim$}spear\\
\glt `(It is) a good plant, we make the spear from it.'
\z

\ea\label{ex:text6-8}
Bea otei tau nao tea mosi koverau,\\
\gll be=a	otei	tau	nao	tea	mosi	koverau\\
     when=\textsc{art}2.\textsc{sg}	man	\textsc{tm}5	go	\textsc{compl}1	cut	koverau.bamboo\\
\glt `When a man is about to cut koverau bamboos,'
\z

\ea\label{ex:text6-9}
toro tara vamataa kurus maa bona maa aba vai\\
\gll toro	tara	va-mataa	kurus	maa	bona	maa	aba	vai\\
     \textsc{tm}7	look	\textsc{advr}-good	very	hither	\textsc{obj}.\textsc{art}2.\textsc{sg}	\textsc{plm}	one2	\textsc{dem}6\\
\glt `(he) must look very closely at the ones'
\z

\ea\label{ex:text6-10}
to pasi mosi maa nae mene o makaree teve.\\
\gll to	pasi	mosi	maa	na=e	mene	o	makaree	teve\\
     \textsc{rel}	\textsc{tm}4	cut	hither	3\textsc{sg}.\textsc{ipfv}=3\textsc{sg}.\textsc{pron}	for	\textsc{art}3.\textsc{sg}	fishing.spear	\textsc{prep}.3\textsc{sg}.\textsc{pron}\\
\glt `that he will cut for his fishing spears.'
\z

\ea\label{ex:text6-11}
Be dee voosu maa bona dee koverau teve,\\
\gll be	dee	voosu	maa	bona	dee	koverau	te-ve\\
     when	carry	go.home	hither	\textsc{obj}.\textsc{art}2.\textsc{sg}	load	koverau.bamboo	\textsc{prep}-3\textsc{sg}.\textsc{pron}\\
\glt `When (he) brings his load of koverau bamboos home,'
\z

\ea\label{ex:text6-12}
ore paa vaasun bona matana inu teori,\\
\gll o=re	paa	vaa-sun	bona	mata-n=a	inu	te=ori\\
     3.\textsc{pron}=\textsc{consec}	\textsc{tm}3	\textsc{caus}-stand	4\textsc{sg}.\textsc{pron}	front-3\textsc{sg}.\textsc{poss}=\textsc{art}2.\textsc{sg}	house	\textsc{prep}=3\textsc{pl}.\textsc{pron}\\
\glt `he stands them up in front of their house,'
\z

\newpage
\ea\label{ex:text6-13}
ei ore roava nana.\\
\gll ei	o=re	roava	nana\\
     \textsc{dem}2	3.\textsc{pron}=\textsc{consec}	in.the.sun	3\textsc{sg}.\textsc{ipfv}\\
\glt `so that they are exposed to the sun.'
\z

\ea\label{ex:text6-14}
Be mahaka, ore paa aeae, ore paa nana mana bari teo suraa.\\
\gll be	mahaka	o=re	paa	aeae	o=re	paa	nana	mana	bari	te=o	suraa\\
     when	dry	3.\textsc{pron}=\textsc{consec}	\textsc{tm}3	light	3.\textsc{pron}=\textsc{consec}	\textsc{tm}3	warm	eventually	4\textsc{pl}.\textsc{pron}	\textsc{prep}=\textsc{art}3.\textsc{sg}	fire\\
\glt `When (they) dry up, they become light, (and) then they (the men) warm them over the fire.'
\z

\ea\label{ex:text6-15}
Pasi nana bata nom ori, beori menamenaa vaatotobin batari bari.\\
\gll pasi	nana	bata	nom	ori	be=ori	mena{$\sim$}menaa	vaa-totobin	bata=ri	bari\\
     \textsc{tm}4	warm	\textsc{cont}	\textsc{ipfv}	3\textsc{pl}.\textsc{pron}	when=3\textsc{pl}.\textsc{pron}	\textsc{red}{$\sim$}bend	\textsc{caus}-straight	\textsc{cont}=3\textsc{pl}.\textsc{ipfv}	4\textsc{pl}.\textsc{pron}\\
\glt `They will warm them, while bending them straight.'
\z

\ea\label{ex:text6-16}
Beve totobin, ore paa kahu bono kosi no koverau,\\
\gll beve	totobin	o=re	paa	kahu	bono	kosi	n=o	koverau\\
     when.3\textsc{sg}.\textsc{pron}	straight	3.\textsc{pron}=\textsc{consec}	\textsc{tm}3	scrape	\textsc{obj}.\textsc{art}2.\textsc{pl}	bottom	3\textsc{sg}.\textsc{poss}=\textsc{art}3.\textsc{sg}	koverau.bamboo\\
\glt `When they (the bamboos) are straight, they scrape the bottom of the koverau bamboos,'
\z

\ea\label{ex:text6-17}
vaamatatopo vo tea vaatopo bono kuna sikiri\\
\gll vaa-matatopo	vo	tea	vaa-topo	bono	kuna~sikiri\\
     \textsc{caus}-be.prepared	\textsc{goal}	\textsc{compl}1	\textsc{caus}-stick	\textsc{obj}.\textsc{art}2.\textsc{pl}	spike.of.fishing.spear\\
\glt `(and) prepare (them) for sticking the spikes (into them).'
\z

\ea\label{ex:text6-18}
Vaatopo vai bona maa kuna sikiri,\\
\gll vaa-topo	vai	bona	maa	kuna~sikiri\\
     \textsc{caus}-stick	now	\textsc{obj}.\textsc{art}2.\textsc{sg}	\textsc{plm}	spike.of.fishing.spear\\
\glt `Now (they) stick in the spikes,'
\z

\ea\label{ex:text6-19}
ore paa pam bari bona kanono.\\
\gll o=re	paa	pam	bari	bona	kanono\\
     3.\textsc{pron}=\textsc{consec}	\textsc{tm}3	fasten.by.tying	4\textsc{pl}.\textsc{pron}	\textsc{obj}.\textsc{art}2.\textsc{sg}	string\\
\glt `and fasten them with a string.'
\z

\ea\label{ex:text6-20}
Be kavara, erau are paa pokaa e bono makaree.\\
\gll be	kavara	erau	a=re	paa	pokaa	e	bono	makaree\\
     when	finished	so	1\textsc{pl}.\textsc{in}.\textsc{pron}=\textsc{consec}	\textsc{tm}3	call	3\textsc{sg}.\textsc{pron}	\textsc{obj}.\textsc{art}3.\textsc{sg}	fishing.spear\\
\glt `When (this is) finished, then we call it the makaree fishing spear.'
\z

\ea\label{ex:text6-21}
Bea abana nao tahii, toro ato momohu tea makaree teori mene tahi vamate bona iana vai a beera to pasi rahi riori tea kanono.\\
\gll be=a	abana	nao	tahii	toro	ato	momohu	te=a	makaree	te=ori	mene	tahi	va-mate	bona	iana	vai	a	beera	to	pasi	rahi	ri=ori	te=a	kanono\\
     when=\textsc{art}2.\textsc{sg}	men	go	sea	\textsc{tm}7	get	first	\textsc{prep}=\textsc{art}3.\textsc{pl}	fishing.spear	\textsc{prep}=3\textsc{pl}.\textsc{pron}	for	spear	\textsc{advr}-dead	\textsc{obj}.\textsc{art}2.\textsc{sg}	fish	\textsc{dem}6	\textsc{art}2.\textsc{sg}	big	\textsc{rel}	\textsc{tm}4	pull	3\textsc{pl}.\textsc{ipfv}=3\textsc{pl}.\textsc{pron}	\textsc{prep}=\textsc{art}2.\textsc{sg}	string\\
\glt `When the men go fishing, (they) must first get their fishing spears for spearing a big fish to death that they will pull (out of the water) with a string.'
\z

\ea\label{ex:text6-22}
O kururu na vaapeha mi nana bono koverau.\\
\gll o	kururu	na	vaa-peha	mi	nana	bono	koverau\\
     \textsc{art}3.\textsc{sg}	kururu.bamboo	\textsc{tm}2	\textsc{caus}-one	with	3\textsc{sg}.\textsc{ipfv}	\textsc{obj}.\textsc{art}3.\textsc{sg}	koverau.bamboo\\
\glt `The kururu bamboo is similar to the koverau bamboo.'
\z

\ea\label{ex:text6-23}
Amaa naono vai to tei me rori bona maa navanava riori: koverau, kururu, koopu, tovu, tabuana.\\
\gll a=maa	naono	vai	to	tei	me	rori	bona	maa	nava{$\sim$}nava	ri=ori	koverau	kururu	koopu	tovu	tabuana\\
     \textsc{art}2.\textsc{sg}=\textsc{plm}	plant	\textsc{dem}6	\textsc{rel}	be.there	with	3\textsc{pl}.\textsc{ipfv}	\textsc{obj}.\textsc{art}2.\textsc{sg}	\textsc{plm}	\textsc{red}{$\sim$}knot	3\textsc{pl}.\textsc{poss}=3\textsc{pl}.\textsc{pron}	koverau.bamboo	kururu.bamboo	koopu.bamboo	tovu.sugarcane	tabuana.sugarcane\\
\glt `The plants that have bars (separated by knots) are: the koverau, kururu and koopu bamboos, the sugar cane and the wild sugar cane.'
\z

\ea\label{ex:text6-24}
O kururu na isuisuvu ruene ni riori.\\
\gll o	kururu	na	isu{$\sim$}isuvu	ruene	ni	ri=ori\\
     \textsc{art}3.\textsc{sg}	kururu	\textsc{tm}2	\textsc{red}{$\sim$}fill.liquid	fresh.water	\textsc{appl}	3\textsc{pl}.\textsc{ipfv}=3\textsc{pl}.\textsc{pron}\\
\glt `The kururu bamboo is used for fetching water.'
\z

\ea\label{ex:text6-25}
Beori tau isuvu ruene ni bari, na dona rori bona kukan a nava kururu.\\
\gll be=ori	tau	isuvu	ruene	ni	bari	na	dona	rori	bona	kukan	a	nava	kururu\\
     when=3\textsc{pl}.\textsc{pron}	\textsc{tm}5	fill.liquid	fresh.water	\textsc{appl}	4\textsc{pl}.\textsc{pron}	\textsc{tm}2	pierce.knots	3\textsc{pl}.\textsc{ipfv}	\textsc{obj}.\textsc{art}2.\textsc{sg}	three	\textsc{art}2.\textsc{sg}	bar	kururu.bamboo\\
\glt `When they going to fetch water with it, (they) pierce (the knots between) three bars.'
\z

\ea\label{ex:text6-26}
Murinae, eori repaa gono bene ta naono eta bebeahu bara kikis\\
\gll muri-nae	e-ori	re=paa	gono	bene	ta	naono	e=ta	bebeahu	bara	kikis\\
     following-3\textsc{sg}.\textsc{poss}.\textsc{pron}	\textsc{h}-3\textsc{pl}.\textsc{pron}	\textsc{consec}=\textsc{tm}3	get	\textsc{obj}.\textsc{art}1.\textsc{sg}	piece	wood	\textsc{art}1.\textsc{sg}=piece	long	and2	strong\\
\glt `For this they get a long and strong stick'
\z

\ea\label{ex:text6-27}
to tahitahi nibusu nao riori bona buaku a nava kururu vaa gina,\\
\gll to	tahi{$\sim$}tahi	nibusu	nao	ri=ori	bona	buaku	a	nava	kururu	vaa	gina\\
     \textsc{rel}	\textsc{red}{$\sim$}pierce	hole	thither	3\textsc{pl}.\textsc{ipfv}=3\textsc{pl}.\textsc{pron}	\textsc{obj}.\textsc{art}2.\textsc{sg}	two	\textsc{obj}.\textsc{art}2.\textsc{sg}	bar	kururu.bamboo	\textsc{lnk}	on.top\\
\glt `with which they pierce holes into the knots of the two top kururu bars,'
\z

\ea\label{ex:text6-28}
beori goe bona nava vaa taeno kururu.\\
\gll be=ori	goe	bona	nava	vaa	tae-n=o	kururu\\
     when=3\textsc{pl}.\textsc{pron}	do.not	\textsc{obj}.\textsc{art}2.\textsc{sg}	bar	\textsc{lnk}	bottom-3\textsc{sg}.\textsc{poss}=\textsc{art}3.\textsc{sg}	kururu.bamboo\\
\glt `while leaving the bar at the bottom of the kururu bamboo intact.'\footnote{lit. `while not doing anything to the bar at the bottom of the kururu bamboo'}
\z

\ea\label{ex:text6-29}
A ruene re goe tea oro tavusu.\\
\gll a	ruene	re	goe	tea	oro	tavusu\\
     \textsc{art}2.\textsc{sg}	fresh.water	\textsc{consec}	do.not	\textsc{compl}1	flow	come.out\\
\glt `Then the water would not flow out.'
\z

\ea\label{ex:text6-30}
O koopu saka pura vamataa haana rapana kasuana ge tea suiri,\\
\gll o	koopu	saka	pura	va-mataa	haa=na	rapa-n=a	kasuana	ge	te=a	suiri\\
     \textsc{art}3.\textsc{sg}	koopu.bamboo	\textsc{neg}1	grow	\textsc{advr}-good	\textsc{neg}2=3\textsc{sg}.\textsc{ipfv}	side-3\textsc{sg}.\textsc{poss}=\textsc{art}2.\textsc{sg}	beach	or	\textsc{prep}=\textsc{art}2.\textsc{sg}	lowlands\\
\glt `The koopu bamboo does not grow well near the beach or in the lowlands,'
\z

\ea\label{ex:text6-31}
na pura vamataa nana tea voora.\\
\gll na	pura	va-mataa	nana	te=a	voora\\
     \textsc{tm}2	grow	\textsc{advr}-good	3\textsc{sg}.\textsc{ipfv}	\textsc{prep}=\textsc{art}2.\textsc{sg}	mountains\\
\glt `(it) grows well in the mountains.'
\z

\ea\label{ex:text6-32}
Meha maa vuri a ta vaan vai to tei rori tea maa vaan vahiava na vaamahaka rori bona tabaan teori komano koopu.\\
\gll meha	maa	vuri	a	ta	vaan	vai	to	tei	rori	te=a	maa	vaan	va-hiava	na	vaa-mahaka	rori	bona	tabaan	te=ori	koma-n=o	koopu\\
     other2.\textsc{sg}	\textsc{plm}	time	\textsc{art}2.\textsc{sg}	people	village	\textsc{dem}6	\textsc{rel}	be.there	3\textsc{pl}.\textsc{ipfv}	\textsc{prep}=\textsc{art}2.\textsc{sg}	\textsc{plm}	village	\textsc{advr}-go.up	\textsc{tm}2	\textsc{caus}-cooked	3\textsc{pl}.\textsc{ipfv}	\textsc{obj}.\textsc{art}2.\textsc{sg}	food	\textsc{prep}=3\textsc{pl}.\textsc{pron}	inside-3\textsc{sg}.\textsc{poss}=\textsc{art}3.\textsc{sg}	koopu.bamboo\\
\glt `Sometimes the people who live in the mountain villages cook their food in the koopu bamboo.'
\z

\ea\label{ex:text6-33}
Na mosi maari bona maa nava koopu,\\
\gll na	mosi	maa=ri	bona	maa	nava	koopu\\
     \textsc{tm}2	cut	hither=3\textsc{pl}.\textsc{ipfv}	\textsc{obj}.\textsc{art}2.\textsc{sg}	\textsc{plm}	bar	koopu.bamboo\\
\glt `(They) cut some bars of koopu bamboo,'
\z

\newpage
\ea\label{ex:text6-34}
eori repaa apeo bona tabaan tabae to rake riori tea apeo: e guu, manii, kaukau, gea tabae to rake riori tea an.\\
\gll e-ori	re=paa	apeo	bona	tabaan	tabae	to	rake	ri=ori	tea	apeo	e	guu	manii	kaukau	ge=a	tabae	to	rake	ri=ori	tea	an\\
     \textsc{h}-3\textsc{pl}.\textsc{pron}	\textsc{consec}=\textsc{tm}3	cook.in.bamboo	\textsc{obj}.\textsc{art}2.\textsc{sg}	food	what	\textsc{rel}	want	3\textsc{pl}.\textsc{ipfv}=3\textsc{pl}.\textsc{pron}	\textsc{compl}1	cook.in.bamboo	\textsc{art}1.\textsc{sg}	pig	possum	sweet.potato	or=\textsc{art}2.\textsc{sg}	what	\textsc{rel}	want	3\textsc{pl}.\textsc{ipfv}=3\textsc{pl}.\textsc{pron}	\textsc{compl}1	eat\\
\glt `and then they cook in the bamboo whatever food they like to cook in the bamboo: pork, possum, sweet potatoes or what they like to eat.'
\z

\ea\label{ex:text6-35}
O koopu vai o peree to apeo riori bona tabaan.\\
\gll o	koopu	vai	o	peree	to	apeo	ri=ori	bona	tabaan\\
     \textsc{art}3.\textsc{sg}	koopu.bamboo	\textsc{dem}6	\textsc{art}3.\textsc{sg}	raw	\textsc{rel}	cook.in.bamboo	3\textsc{pl}.\textsc{ipfv}=3\textsc{pl}.\textsc{pron}	\textsc{obj}.\textsc{art}2.\textsc{sg}	food\\
\glt `(It is) the koopu green bamboo in which they cook the food.'
\z

\ea\label{ex:text6-36}
Eori repaa vaamahaka bari teo suraa.\\
\gll e-ori	re=paa	vaa-mahaka	bari	te=o	suraa\\
     \textsc{h}-3\textsc{pl}.\textsc{pron}	\textsc{consec}=\textsc{tm}3	\textsc{caus}-simmer	4\textsc{pl}.\textsc{pron}	\textsc{prep}=\textsc{art}3.\textsc{sg}	fire\\
\glt `Then they let it simmer on the fire.'
\z

\ea\label{ex:text6-37}
Beori vaatoka bono koopu teo suraa,\\
\gll be=ori	vaa-toka	bono	koopu	te=o	suraa\\
     when=3\textsc{pl}.\textsc{pron}	\textsc{caus}-perch	\textsc{obj}.\textsc{art}3.\textsc{sg}	koopu.bamboo	\textsc{prep}=\textsc{art}3.\textsc{sg}	fire\\
\glt `When they put the koopu bamboo onto the fire,'
\z

\ea\label{ex:text6-38}
toro vavereverete bata ni rori bono koopu to mene asi.\\
\gll toro	va-vere{$\sim$}verete	bata	ni	rori	bono	koopu	to	mene	asi\\
     \textsc{tm}7	\textsc{mult}-\textsc{red}{$\sim$}turn	\textsc{cont}	\textsc{appl}	3\textsc{pl}.\textsc{ipfv}	\textsc{obj}.\textsc{art}3.\textsc{sg}	koopu.bamboo	\textsc{rel}	\textsc{tm}6	burn\\
\glt `(they) must repeatedly turn the koopu bamboo around so that it won't burn.'\footnote{lit. `... the bamboo that otherwise would burn.'}
\z

\newpage
\ea\label{ex:text6-39}
Eara na von rara a koopu tea ta vahiava mene paku bono takinima na maa inu teara.\\
\gll e-ara	na	voni	rara	a	koopu	te=a	ta	vahiava	mene	paku	bono	takinima	n=a	maa	inu	te=ara\\
     \textsc{h}-1\textsc{pl}.\textsc{in}.\textsc{pron}	\textsc{tm}2	buy	1\textsc{pl}.\textsc{in}.\textsc{ipfv}	\textsc{art}3.\textsc{pl}	koopu.bamboo	\textsc{prep}=\textsc{art}2.\textsc{sg}	people	inland	for	make	\textsc{obj}.\textsc{art}2.\textsc{pl}	bamboo.wall	3\textsc{sg}.\textsc{poss}=\textsc{art}2.\textsc{sg}	\textsc{plm}	house	\textsc{prep}=1\textsc{pl}.\textsc{in}.\textsc{pron}\\
\glt `We buy the koopu bamboo from the mountain people to make the bamboo walls of our houses.'
\z

\ea\label{ex:text6-40}
Beara taboo a inu teara bono takinima ore paa tara vamataa.\\
\gll be=ara	taboo	a	inu	te=ara	bono	takinima	o=re	paa	tara	va-mataa\\
     when=1\textsc{pl}.\textsc{in}.\textsc{pron}	wall	\textsc{art}2.\textsc{sg}	house	\textsc{prep}=1\textsc{pl}.\textsc{in}.\textsc{pron}	\textsc{obj}.\textsc{art}2.\textsc{pl}	bamboo.wall	3.\textsc{pron}=\textsc{consec}	\textsc{tm}3	look	\textsc{advr}-good\\
\glt `When we wall our house with bamboo walls,  it looks beautiful.'
\z
